\section{International climate policy: majority support even when people aware of the costs}

% AF: premier paragraphe pas très informatif, on ne comprend pas ce que font les papiers cités
The preceding section showed that citizens may accept supranational coordination when it is democratic, credible, and confined to genuinely global issues. A separate objection concerns the private and fiscal costs of the policies that such institutions would implement. Support for international redistribution and climate action may be overstated when respondents do not understand the implications for household income, energy prices, or domestic public budgets. The analysis must therefore move from abstract approval to the evaluation of explicitly costly policies. The relevant question is not simply whether citizens bear a cost, but how that cost is distributed, whether the policy is expected to reduce emissions effectively, and whether other countries contribute to the collective effort. Each dimension can affect acceptance, although broad international participation is not always necessary for substantial support \citep{bechtel_mass_2013,bernauer_how_2015,mcgrath_how_2017}.

Combining climate mitigation with redistribution is not merely a matter of political framing. Without compensatory transfers, ambitious mitigation policies may increase extreme poverty, whereas progressive recycling of carbon-pricing revenues and additional international climate finance can offset these regressive effects \citep{soergel_combining_2021}. The Global Climate Scheme (GCS) evaluated by \citet{fabre_public_2026} provides a demanding test of the political acceptability of such an integrated policy. Under the GCS, revenues from a global cap-and-trade system are distributed equally across adults, lifting 600 million people out of extreme poverty while imposing an average net cost of \$90 per month in the United States and \texteuro45 per month in Germany. The International Climate Scheme (ICS) applies the same principle under partial international participation. Respondents evaluate Low-, Mid-, and High-coverage scenarios accounting for 33\%, 56\%, and 72\% of global emissions, respectively, as well as a fourth variant displaying country-level distributive effects 
in the High-coverage scenario. % AF: added

% AF: il faut d'abord citer le soutien au GCS dans le papier NHB
Despite its explicit costs, the GCS is accepted by 55\% of respondents, with country-level acceptance ranging from 49\% in the United States and Russia to 85\% in Saudi Arabia \citep{fabre_public_2026}. Acceptance of the ICS increases only modestly with country coverage, from 65\% in the Low scenario to 68\% in the High scenario, a difference of 4 %3.8 AF: changed
percentage points. Among European respondents, acceptance rises when China joins the coalition but does not increase further when other high-income countries participate. By contrast, respondents in the United States and Japan react more strongly % AF: more strongly or only?
to the combined participation of China, the European Union, and other high-income countries. The effect of coverage is entirely driven by the 74\% of respondents who correctly understand that the policy would increase gasoline prices. Even under low coverage, respondents appear willing to bear the ICS's higher cost in exchange for guaranteed poverty alleviation and decarbonization in the Global South. Its greater acceptance relative to the GCS may reflect the greater credibility of a proposal that explicitly identifies participating countries, although the evidence does not isolate this mechanism.

% Overall, the evidence does not support a mechanical rule under which majority acceptance follows whenever two of three objectives---fairness, environmental effectiveness, and individual self-interest---are satisfied. These dimensions instead interact. % AF: je ne comprends pas pourquoi tu écris ça : c'est justifié par quels résultats ? Dans Douenne & Fabre (22), on trouve un soutien majoritaire dès qu'on a deux perceptions sur trois. Peut-être que d'autres trouvent autre chose, mais dans ce cas-là il faut citer et justifier
% AF: faut citer de façon plus précise les papiers ci-dessous
Higher individual costs generally reduce support, whereas equitable burden-sharing, broader participation, and credible enforcement increase the perceived legitimacy and effectiveness of international climate agreements \citep{bechtel_mass_2013}. In the specific context of North--South climate finance, acceptance is also higher when other donor countries contribute, allocation decisions % AF: comment ça allocation decisions? peu clair
are shared between donors and recipients, and recipient institutions are perceived as effective \citep{gampfer_obtaining_2014}. International reciprocity is not, however, an absolute prerequisite: substantial support for unilateral climate action persists even when other countries do not adopt equivalent policies \citep{bernauer_how_2015,tingley_conditional_2014,mcgrath_how_2017}. Perceived fairness and environmental effectiveness may therefore partly offset the influence of narrow material self-interest. The political challenge is not to identify a ``perfect'' policy, but to design a credible and transparent scheme whose environmental benefits, redistributive consequences, and burden-sharing arrangements are clearly understood.

These findings establish that costly international climate policies can attract conditional majority acceptance. The remaining question is whether this acceptance extends to a willingness to contribute a clearly specified share of personal income.


